\section{工程可视化与监控}

\subsection{Dashboard 设计}
Compound AI 的每个模块都要在运营面板上有可见的仪表盘：trace fidelity curve 做横向对比，judge reward curve 与 user satisfaction signal 汇总成二合一视图，solver output 则在另一个 panel 中突出 constraint violation rate。下表列出关键仪表盘与推荐的刷新频率。
\begin{table}[H]
\centering
\begin{tabular}{lp{8cm}r}
\toprule
仪表盘 & 包含指标 & 建议刷新频率 \\
\midrule
Trace health & fidelity、cost variance、skew coverage & 1 分钟 \\
Judge stability & accuracy/proactivity/credibility 分布、NPS proxy scores & 30 秒 \\
Constraint guard & solver violation rate、brush 工艺 compliance & 10 秒 \\
\bottomrule
\end{tabular}
\caption{Compound AI 主要 dashboard 与刷新策略}
\end{table}
\begin{knowledgebox}{可视化的反馈节奏}
把不同模块的监控刷新频率按 criticality 分层，避免同一时刻刷新所有指标导致监控卡顿；同时用 \texttt{trace\_id} 作为通用维度，在 dashboard 上串联 trace、judge、solver 三段信息，帮助工程师跨模块排查。
\end{knowledgebox}

\subsection{Trace 与 Plan 回放}
真正的透明度来自 trace/plan 的回放能力：每个 plan 执行后都写入 \texttt{trace\_id.json}，包含 heuristic 得分、约束违例和 judge reward，从而可以在运维网页里逐步播放 search trace、judge dialogue、solver output。为了让每次回放都有视觉印象，我们把关键 frame 的时间戳写入同一条 log，便于回放时加载对应的 slide 截图或 frame 图像。
\begin{knowledgebox}{回放提升审计速度}
一条 trace 回放不仅能还原 decision path，还能让 engineer 直接看见与 slide 对齐的 visual evidence。比如 00:16:08--00:16:22 段的 APAC slide，会在回放里作为 judge trigger 的注释；00:32:25--00:32:45 段的 DSPy 架构图则用来解释 solver 的 step-by-step 计算。
\end{knowledgebox}
\begin{warningbox}{回放数据不可缺}
缺失 trace/frame 之间的时间对齐会拖慢调查进度。务必在日志里同时写入 \texttt{trace\_id}、\texttt{frame\_path}、\texttt{timestamp}，否则回放时只能靠模糊对照，运营判断会失真。
\end{warningbox}

\subsection{本章小结}
工程监控的重点是把 metrics、visual evidence、回放逻辑合并为一个可检索的三链路，让事件发生时可以迅速定位到对应的 trace frame、judge dialog 和 solver plan，从而对故障做出精准响应。

